To analyze long-term structural changes, we compare the five-year periods 2001–2005 and 2021–2025. This timeframe smooths out annual fluctuations and, through a broader data base, ensures the statistical validity of the results.
The analysis is conducted on two levels: 1. Recording performance trends and density for all groups to identify group differences. 2. Examining the development to identify the structure in which performance changes occur.
This allows us to identify general trends and determine the structure of the changes.

To measure the performance density, we define the performance gap ($G_p$) between the three best and the three worst (e.g., 33rd–35th place) in our data set for the two time periods. The averaging of the top 3 reflects the international standard of success (podium), while the bottom 3 provide a stable and comparable baseline. We used the performance gap to measure the breadth of performance at the top level and utilized all available information (e.g., rankings) from our dataset to measure polarization. Mathematically, the gap ($G_p$) is defined as follows:

\begin{equation}
G_p = \left( \frac{1}{3|P|} \sum_{y \in P} \sum_{r=1}^{3} x_{y,r} \right) - \left( \frac{1}{3|P|}  \sum_{y \in P} \sum_{r=33}^{35} x_{y,r} \right)
\end{equation}

where $x_{y,r}$ represents the individual performance of rank $r$ in year $y$, and $|P|$, in our case $|P|=5$, represents the number of years per period. The development of the performance gap was calculated as the relaive percentage change between the two five-year periods.

To analyze the development of the different rankings, we calculate the arithmetic mean for each rank across all groups for the start and end periods. For each rank, even when differentiating by age group and gender, the sample size is at least n = 80. According to the central limit theorem, we can assume a normal distribution and thus calculate the arithmetic mean. The development of each rank was determined by calculating the percentage change between the start and end periods, normalized to the start period. A positive value indicates an improvement in the average ranking, a negative value a decline.

During the period under review, several rule changes were implemented regarding hurdle heights and throwing weights, particularly in the U18 (2012) and U20 age groups (2001) \citep{ProgressionWorldAthletics}. Therefore, the affected disciplines were excluded from median and rank analyses, as changes in these parameters are directly attributable to equipment changes and not to athletic development. However, the performance gap analysis was retained, as this variance can be calculated independently of the absolute values.